\cotefile{135}
\scanpage{3}
``Tordre'' une catégorie de Picard

1) Multiplier une Gr-cat par un groupe

Soit $C$ une (gr)-cat, $\Gamma$ un groupe, on va définir une nouvelle (Gr)-cat. notée $\Gamma\times C$. Le plus simple est de dire que $\Gamma$ désigne une catégorie discrète d'ens. sous-jacent $\Gamma$, et la loi de composition de $\Gamma$ définit un foncteur $\otimes$, avec une unique \deleted{structure}\inserted{contrainte} d'associativité. \deleted{\unsure{manière|Ceci|évidente}} [Il existe une \deleted{contrainte}\deleted{\unsure{contrainte}}\inserted{contrainte} de commutativité ssi $\Gamma$ est commutatif. (Les deux contraintes sont alors compatibles, et $\Gamma$ est une Catégorie \unsure{Abl|Abél|Ab|ACU})] On prend alors le produit des (Gr-cats) $\Gamma$ et $C$.
\[
\begin{aligned}
\pi_0(\Gamma\times C)&=\pi_0(\Gamma)\times\pi_0(C)=\Gamma\times\pi_0(C)\xrightarrow{\unsure{p|\pi}=\text{projection}}\pi_0(C)\\
\pi_1(\Gamma\times C)&=\pi_1(\Gamma)\times\pi_1(C)=\pi_1(C)\\
\unsure{k|K}_{\Gamma\times C}&=p^*(k_C)
\end{aligned}
\]

2) Supposons que $C$ soit $=$ une Pic-Cat, et \deleted{\unsure{une}}\inserted{$\Gamma$} commutatif. \deleted{Alors} $\Gamma\times C$ est une Pic-catégorie.

3) Si $C$ est une Pic-catégorie, les \inserted{autres} \deleted{structures}\inserted{contraintes} de comm. qui en font (avec sa \deleted{\unsure{contrainte}} d'associativité) une Pic-catégorie -- \unsure{i.e.}\deleted{\unsure{t.q.|types}} \deleted{\unsure{sont}} compatibles avec la contrainte d'ass., correspondent aux applications \uline{bilinéaires} \uline{antisymétriques}
\[
\varphi\colon\pi_0(C)\times\pi_0(C)\longrightarrow\pi_1(C)
\]
L'application correspondante
\[
\pi_0(C)\xrightarrow[\Sigma_\varphi]{}\pi_1(C)
\]
est donnée par
\[
\Sigma_\varphi(x)=\varphi(x,x)
\]
% archive sheet number: 2
\scanpage{4}
\begin{center}
LONELY ROAD \qquad (a song)
\end{center}
\noindent\begin{minipage}[t]{.48\linewidth}\small\raggedright
You \enspace brother\\
\deleted{\unsure{b|h}}and me\\
and all of us\\
brothers and sisters\\
each one is going his on\\
his own\\
lonely road

\medskip
Some will meet\\
and some will part\\
some go alongside\\
a little while\\
or a long while\\
or a life's while -\\
but always remember\\
oh brother !\\
that you and me\\
and each and everyone\\
has his own lonely road -\\
his only road\\
that's only his

\medskip
Thousands will look thesame\\
millions will look the same\\
they'll go under just one name -

\smallskip
but just look closer\\
brother\\
and you will see\\
none is like the other\\
no two are thesame\\
each one of the millions\\
even as his name\\
has got his own\\
his single\\
lonely road
\end{minipage}\hfill
\begin{minipage}[t]{.48\linewidth}\small\raggedright
Your road\\
brother\\
is not that concrete highway\\
that thousands take -

\smallskip
it's just the path\\
that you yourself are treading\\
through the high grasses\\
of the grasslands\\
through the shrubs of the forests\\
the wandering \unsure{sánds|sands} of the deserts\\
the rolling waves of the seas\\
of the seas
\end{minipage}\par
% archive sheet number: 3
\scanpage{5}
4) \uline{Lemme} Soient $M,N$ des groupes commutatifs, et $\Sigma\colon M\longrightarrow {}_2N$ un hom. Alors il existe une application bilinéaire antisymétrique $\varphi\colon M\times M\longrightarrow N$ telle que $\varphi(x,x)=\Sigma(x)$ pour tt $x\in M$.

Remplaçant $M,N$ resp. par $\deleted{\unsure{\Sigma|2}}M_2,{}_2N$, on peut supposer $M,N$ des vectoriels \deleted{\unsure{de dim}} sur $\mathbb F_2$.

Soit $(e_i)_{i\in I}$ une base de $M$, \deleted{\unsure{on}} Alors $\deleted{\varphi}\inserted{\Sigma}$ est donnée par les $n_i=\Sigma(e_i)\in N$, tandis que $\varphi$ sera donnée par les $\varphi(e_i,e_j)=n_{ij}\in N$; la condition d'antisymétrie \inserted{(= symétrie)} est que $\deleted{\unsure{\varphi(e_i,e_j)=}} n_{ij}=n_{ji}$. On prendra $\varphi(e_i,\unsure{e_j|j})=\begin{cases}0&\text{si }i\ne j\\ n_i&\text{si }i=j\end{cases}$ \unsure{O.K.|OK}

\inserted{\uline{Corollaire} Toute catégorie de Picard devient \uline{stricte} en changeant sa seule contrainte de commutativité. En tant que Gr-catégorie, elle est déjà triviale\ldots}

5) Soit $C$ une catégorie de Picard, $\Gamma$ un groupe commutatif, \deleted{$\varphi$}
\[
\varphi\colon\Gamma\times\Gamma\longrightarrow\pi_1(C)
\]
une application bilinéaire antisymétrique. On désigne par $\Gamma\overset{\varphi}{\times}C$ la catégorie de Picard déduite de $\Gamma\times C$ (cf n° 2) par l'application bilinéaire \unsure{antisym|antisym.}
\[
\bar\varphi\colon\pi_0(\Gamma\times C)\times\pi_0(\Gamma\times C)\longrightarrow\pi_1(\Gamma\times C)
\]
(cf n° 3) donnée par
\[
\bar\varphi\big((\alpha,x),(\beta,y)\big)=\varphi(\alpha,\beta)
\]
i.e. composée \deleted{\unsure{de}}
\diagram{\begin{tikzcd}[ampersand replacement=\&,column sep=small]
{\pi_0(\Gamma\times C)\times\pi_0(\Gamma\times C)} \arrow[dr] \& \& {\pi_1(\Gamma\times C)\simeq\pi_1(C)}\\
\& {\deleted{\unsure{\pi_0(}}\Gamma\times\Gamma} \arrow[ur,bend left=10,draw=none,"\deleted{\longrightarrow}"{description,sloped}] \arrow[ur,bend right=7,"\varphi"'] \&
\end{tikzcd}}

6)
% archive sheet number: 4
\scanpage{6}
\begin{center}
\deleted{LONELY ROAD \qquad (a song)}
\end{center}
\noindent\begin{minipage}[t]{.48\linewidth}\small\raggedright
\deleted{You \enspace brother\\
\deleted{\unsure{b|h}}and me\\
and all of us\\
brothers and sisters\\
each one is going his on\\
his own\\
lonely road}

\medskip
\deleted{Some will meet\\
and some will part\\
some go alongside\\
a little while\\
or a long while\\
or a life's while -}

\smallskip
\deleted{but always remember\\
oh brother !\\
that you and me\\
and each and everyone\\
has his own lonely road -\\
his only road\\
that's only his}

\medskip
\deleted{Thousands will look thesame\\
millions will look the same\\
they'll go under just one name -}

\smallskip
\deleted{but just look closer\\
brother\\
and you will see\\
none is like the other\\
no two are thesame\\
each one of the millions\\
even as his name\\
has got his own\\
his single\\
lonely road}
\end{minipage}\hfill
\begin{minipage}[t]{.48\linewidth}\small\raggedright
\deleted{Your road\\
brother\\
is not that concrete highway\\
that thousands take -}

\smallskip
\deleted{it's just the path\\
that you yourself are treading\\
through the high grasses\\
of the grasslands\\
through the shrubs of the forests\\
the wandering \unsure{sánds|sands} of the deserts\\
the rolling waves of the seas\\
of the seas}
\end{minipage}\par
% archive sheet number: 5
\scanpage{8}
\mspage{1}
1. $K$ complexe semi-simplicial

$\unsure{\mathscr P|P}$ \deleted{Catégorie} $\otimes$ Catégorie AU

représentation \inserted{$\Lambda$} de $K$ dans $\mathscr P$
\[
\text{données}\left\{
\begin{aligned}
&\text{a) }\forall u\in K_1,\qquad\Lambda(u)\in\mathscr P\\
&\text{b) }\forall\sigma\in K_2,\ \text{posant }u=d_1\sigma,\ v=d_3(\sigma),\ w=d_{\unsure{2}}(\sigma)\\
&\hspace{9em}\Lambda(u)\otimes\Lambda(w)\simeq\Lambda(\unsure{v|u|w})
\end{aligned}\right.
\]
axiome \quad c) $\forall\Delta\in K_3$, on a ``axiome du pentagone

La composée \quad $K_1\longrightarrow\mathscr P\longrightarrow\pi_0(\mathscr P)$ définit un ``homo\uline{morph}isme'' \quad $K_1\xrightarrow{\alpha}\pi_0(\mathscr P)=G$

d'où un ``homomorphisme'' \quad $K_1\xrightarrow{\bar\alpha}\operatorname{Aut}\big(\underbrace{\pi_1(\mathscr P)}_{M}\big)$

d'où un système local sur $K$ de valeur $\underline M$, noté $(\unsure{{}_\alpha M|M},\bar\alpha)=\bar M$

\uline{Pb} Quand $K,\mathscr P,\deleted{\unsure{(M,\bar\alpha)}}\inserted{\alpha\ \text{(donc }\bar M\text{)}}$ sont fixés, on veut ``classer'' les représentations de $K$ \unsure{dans|ds} $\mathscr P$ correspondantes.

% slanted
\margin{$\operatorname{REP}(K;\mathscr P,\alpha)$\\$\operatorname{REP}(\mathcal X;\mathscr P,\alpha)$}
\uline{Résultat} Ces représentations forment une 2-catégorie, qui est un ``2-torseur'' sous une \uline{2-catégorie de Picard} \inserted{stricte} d'invariants
\[
H^2(K,\unsure{(}\bar M\deleted{\unsure{,\bar\alpha)|,\alpha)}}),\quad
H^1(K,\unsure{(}\bar M\deleted{\unsure{,\bar\alpha)|,\alpha)}}),\quad
H^0(K,\unsure{(}\bar M\deleted{\unsure{,\bar\alpha)|,\alpha)}}).
\]
Cette 2-catégorie est contravariante p.r. à $K$ variable, \deleted{\illegible{2}} covariante p.r. à $\mathscr P$ variable

Cas particulier utilisé par la suite : $K$ est le \uline{nerf} d'une catégorie $\mathcal X$. \deleted{\unsure{On}}

2. $\mathscr P$ Gr-catégorie \deleted{\unsure{de Picard}} donnée, d'invariants
\[
\pi_0(\mathscr P)=\unsure{C|G}\ ;\quad\pi_1(\mathscr P)=M
\]
On va lui associer une 2-gerbe \inserted{$\widetilde{\mathscr P}$} sur $B_G$, donc à tout $X\in\operatorname{ob}B_G$ (ensemble où $\unsure{G|C}$ opère) on va associer une catégorie $\widetilde{\mathscr P}(X)$, ainsi :

Soit $\mathcal X$ le nerf de la catégorie \inserted{$\mathcal X$} sous-jacente à $X$
% archive sheet number: 7
\scanpage{9}
($\operatorname{Ob}\mathcal X=X,\ \operatorname{Fl}\mathcal X=G\times X\ldots$), \deleted{\unsure{on prend}}\inserted{$\alpha_X$ l'}application canonique
\[
K_1=\operatorname{Fl}\mathcal X=\unsure{G|C}\times X\xrightarrow{\mathrm{pr}_1}G
\]
On pose
\[
\widetilde{\unsure{\mathscr P|P}}(X)=\operatorname{Rep}(\mathcal X;\mathscr P,\alpha_X)=\operatorname{Rep}(X,\mathscr P,\alpha_{\unsure{x|X}})
\]
avec la contravariance p.r. à $X$ déjà vue. On trouve que c'est un 2-champ, qui est $=$ une 2-gerbe liée par $M_{(B_G)}$ \deleted{\unsure{(muni de|construit}} (défini par l'op. de $G$ sur $M$).

\deleted{\unsure{3.}} \uline{\unsure{Variance|Varions|Variante}} $\widetilde{\mathscr P}$ est covariant en $\mathscr P$ dans un sens évident \ldots\ \deleted{\unsure{On a un fon}}\inserted{En fait \unsure{un|-}} un morphisme de la 3-catégorie des \deleted{3. Fixons} $\mathscr P$ dans la 3-catégorie des 2-gerbes sur des $B_G$ variables.

3. On a ainsi une équivalence de 3-catégories. Plus \unsure{spécifiquement|explicitement}, $G$ et le $G$-module $M$ étant fixés, on a une équivalence entre la 3-catégorie des Gr-catégories \deleted{de Picard} épinglées par $(G,M)$, et la 3-catégorie des 2-gerbes sur $B_G$ liées par $M_{(B_G)}$

Soient
\[
\begin{aligned}
&\mathscr P &&\text{Gr-catégorie }\deleted{\unsure{\text{définie}|\text{d'invariants}}}\inserted{\text{épinglée par}} &&G,M\\
&\mathscr P' &&\rule{5em}{.3pt} &&G',M'
\end{aligned}
\]
\[
G\overset{u}{\unsure{\relbar|\longrightarrow}}G',\quad M\overset{\unsure{v|o}}{\relbar}M'\qquad\text{hom compatibles aux opérations}
\]
Alors les $\otimes$ hom $\mathscr P\relbar\mathscr P'$ compatibles avec $G\relbar G', M\relbar M'$ forment une \uline{2-catégorie}
% archive sheet number: 8
\scanpage{10}
\mspage{2}
qui est un 2-\inserted{pseudo-}torseur sous une 2-Pic-catégorie, et ce 2-torseur est 2-équivalent à celui des hom de 2-gerbes sur $B_G$ liées par $M'_{B_G}$ de
\margin{$M_{B_G}\unsure{\relbar|\longrightarrow}M'_{B_G}$\\[3pt]$B_G\xrightarrow{B(u)}B_G$}
$\widetilde{\unsure{\mathscr P|P}}\otimes_{M_{B_G}}M'_{B_G}$ dans $B_{\unsure{u}}^*(\widetilde{\mathscr P}')$. Les invariants de la 2-Pic-catégorie pertinente sont $H^2(G,M'),H^1(G,M'),H^0(\unsure{C|G},M')$. L'obstruction à trouver un objet dans le 2-ps.-torseur est dans $H^3(G,M')$, et cela se calcule comme différence de deux classes, images canoniques des éléments qu'on devine dans $H^3(C,M)$ et $H^3(C',M')$.

4. Généralisation aux Gr-champs sur un topos $S$. \deleted{\unsure{On les|Bien sûr}} Si $G$ (groupe sur $S$) et $M$ ($G$-Module sur $S$) sont donnés, on interprète tout en termes de 2-gerbes ``relatives'' sur $B_G$ mod $S$ liées par $M_{B_G}$, i.e. 2-gerbes sur $B_G$ trivialisées sur $S$ (par $S\simeq B_e\longrightarrow B_G$)

Il faudrait le mettre ``en \unsure{série|séries|scène}'' avec l'histoire \unsure{des|de} \uline{1-gerbes} \unsure{relative|relatives} sur $B_C$, correspondant aux \uline{extensions} de $G$ par $M$. Il faudrait traiter en $=$ temps le cas d'un $M$ non commutatif\ldots
% archive sheet number: 9
\scanpage{11}
5. Relation avec les ``noyaux'' de Mac-Lane, ou ``liens''. Soit $\unsure{G|C}$ groupe, $F$ groupe,
\[
\varphi\colon G\longrightarrow\operatorname{Aut\,ext}(F)
\]
hom (\unsure{rev. à|revient à|cela $=$} se donner un lien sur $B_G$). On va lui associer une \inserted{Gr-}catégorie \deleted{de Picard} $\unsure{\mathscr P|P}(\varphi)$ d'invariants $G$ et $\deleted{\unsure{M}}\ M=\mathrm Z(F)$ (centre)

a) \uline{Description directe}. On a la \inserted{Gr} catégorie $\mathscr P(F)$ des bitorseurs sous $F$, qui a comme invariants
\[
\begin{aligned}
\pi_0(\mathscr P(F))&\in\operatorname{Aut\,ext}F\\
\pi_1(\mathscr P(F))&=\mathrm Z(F)\qquad\text{centre de }F
\end{aligned}
\]
On définit donc $\mathscr P(\varphi)$ comme image inverse.

\deleted{b) Description en termes des 2-gerbes liées par $M$.}

\uline{NB} Les trivialisations de $\mathscr P(\varphi)$ correspondent aux réalisations de $\varphi$ par une extension de $G$ par $F$ subordonnée à $\varphi$. Il faudrait encore préciser ce point en termes \unsure{de|des} l'équivalence de 2-catégories (2-torseurs sous une 2-catégorie de Picard). On \unsure{retrouve|retrouvera} ainsi la classe de Mac-Lane dans $H^3(G,M)$, obstruction à la réalisabilité du ``noyau'' $\varphi$ par une extension (ou encore du lien $\varphi$ par une 1-gerbe). \deleted{\unsure{Notons}}\inserted{Mais} \deleted{\unsure{remarquons}}\inserted{le fait que} \deleted{\unsure{tous}} \unsure{élément|élém\textsuperscript{t}|éléments} de $H^3(\unsure{C|G},M)$ (i.e. tte catégorie de Picard \unsure{épinglée|épinglées} par $(G,M)$) puisse ainsi s'obtenir, à l'aide d'un ``noyau'', ne sort pas \inserted{(pour l'instant)} \unsure{d'un fait général|d'un résultat général|d'une logique générale}.
% archive sheet number: 10
\scanpage{12}
\mspage{3}
c) Soient $\unsure{\mathscr P|P},\mathscr P'$ deux catégories de Picard, épinglées par $M,N$\unsure{.|!} Pour qu'ils soient \unsure{équivalentes|équivalents} (avec leur épinglage) il f. et s. que les applications $\sigma_{\mathscr P},\sigma_{\mathscr P'}\colon M\longrightarrow {}_2N$ soient les \unsure{mêmes|$\overline{\mathrm{m}}$s|$=$s}. La catégorie des équivalences est alors \deleted{une}\inserted{un 1-torseur sous une} catégorie de Picard \uline{stricte}
% slanted
\inserted{d'invariants $\operatorname{Ext}^{1}(M,N)$ et $\operatorname{Hom}(M,N)$ \unsure{(définie par|(obtenus par} $R\underline{\operatorname{Hom}}(M,N)$ tronqué\unsure{\ldots})}

d) La classification à équivalence près des Cat. de Picard épinglées par $M,N$ est (comme groupe) $\simeq\operatorname{Hom}(M,{}_2N)$.

e) Revenant à la situation de c) et utilisant la multiplication de Baer des Pic cat. épinglées par $(M,N)$, on trouve que \deleted{\unsure{celles}} celles qui correspondent à un $\sigma\colon M\longrightarrow{}_2N$ donné forment une 2-catégorie qui est un \unsure{2|1}-torseur sous la 2-catégorie de Picard stricte formée des Cat de Pic. \uline{strictes} épinglées par $M,N$ \deleted{\unsure{E}} [qui est un 2-groupoïde connexe : Tous les 0-objets sont 2-équivalents,\ldots] \deleted{\unsure{les}}

f) Les Pic-catégories forment une \uline{2-$\otimes$-catégorie \deleted{\unsure{de}} AUC} de façon évidente, \deleted{\unsure{Si $\mathscr P,\mathscr P'$ et $\mathscr P$}} (on s'envoie dans \inserted{(et $=$ \uline{sur} les objets et flèches)} la 1-$\otimes$-catégorie \inserted{AUC} sous-jacente à la catégorie \deleted{\unsure{des ob}}\inserted{$D$}
% archive sheet number: 11
\scanpage{13}
6. Pic-catégories.

On peut définir les Pic-catégories comme éléments d'une 3-catégorie, qui soit une sous-3-catégorie (pleine \inserted{\deleted{\unsure{à}} au niveau des 2-objet et 3-objet)} de celle des Gr-catégories. \deleted{\unsure{Ainsi|Et si|Soient} $\unsure{si}\ \unsure{\mathscr P|P},\mathscr P'$ sont deux Pic-catégories, et $\Lambda,\Lambda'\colon\mathscr P\rightrightarrows\mathscr P'$ deux $\otimes$ foncteurs \unsure{AC|AU}} Mais on peut aussi \unsure{penser|passer} à la 2-catégorie associée en gardant les $=\mathrm s$ 0 et 1 objets, mais en prenant les 2-objets à isom. près \inserted{(c'est en fait la relation triviale)}; \deleted{\unsure{on}} et \inserted{oubliant les 3-objets} se restreint aux cat. de Picard strictes, on retrouve alors la 2-catégorie de \uline{Deligne} (où les catégories \uline{Hom} sont elles-m$\overline{\rm m}$ des catégories de Picard déduites des $R\operatorname{Hom}(L^{\displaystyle\cdot},L'^{\displaystyle\cdot})$ par troncature).

Si on n'insiste pas sur les Pic-cat. \uline{strictes}, on peut dire ceci :

a) Soit $\mathscr P$ une Gr-catégorie. Elle admet une structure de Pic-catégorie compatible avec \unsure{le reste|la vente|la préd\textsuperscript{te}} (i.e. une \deleted{structure} contrainte de comm. compatible avec \unsure{le reste|la contr.}) ssi elle est splittable, i.e. $k(\mathscr P)\in H^3(\overset{\substack{M\\ \Vert}}{\pi_0},\overset{\substack{N\\ \Vert}}{\pi_1})$ est nulle Alors $\mathscr P$ admet $=$ une structure de Picard stricte.

b) Les structures de Picard \deleted{strictes} sur $\mathscr P$ forment un torseur sous $\operatorname{Bil}\deleted{\unsure{\mathrm{alt}|\mathrm{ant}}}\unsure{\mathrm{A}|\mathrm{Alt}|\mathrm{ant}|\mathrm{alt}}(\pi_0,\pi_0;\pi_1)$. Les structures de Picard strictes forment un sous-torseur sous $\operatorname{Bil\,alt}(\pi_0,\pi_0;\pi_1)$.
% archive sheet number: 12
\scanpage{14}
\mspage{4}
\deleted{\unsure{Pic}} additive ordinaire formée des triples $(M,N,\sigma\colon M\longrightarrow{}_2N)$. Si $\unsure{\mathscr P|P},\mathscr P'$ sont associés à $(M,N,\sigma),(M',N',\sigma')$, alors deux foncteurs $\Lambda,\Lambda'\colon\mathscr P\longrightarrow\mathscr P'$ sont isom. ssi. ils définissent le $=$ hom $(M,N,\sigma)\longrightarrow(M',N',\sigma)$ (donc $D$ ``est'' la catégorie \unsure{déduite|discrète} de $(\mathrm{Pic})$ en passant au quotient par les 2-flèches).
% slanted
\margin{NB $\underline{\operatorname{Hom}}_{\otimes\,\mathrm{AUC}}(\mathscr P,\mathscr P')$ est un groupoïde}
Mais bien sûr, la catégorie des $\otimes$-foncteur AUC $\mathscr P\longrightarrow\mathscr P'$ qui sont subordonnés à un hom donné $(M,N,\sigma)\mathbin{\unsure{-|\to}}(M',N',\sigma)$ forment un groupoïde qui est un torseur sous la Pic-catégorie définie par \unsure{tronsection|transection|troncature|troncation} à partir de $R\underline{\operatorname{Hom}}(M,N')$ (d'invariant $\operatorname{Ext}^{\unsure{1}}(M,N'),\operatorname{Hom}(M,N')$.

7) \deleted{\unsure{Présentation|Il résulte|Traductions}} \uline{Extension aux champs de Picard sur un topos}

Ils forment une 2-\inserted{$\otimes$-}catégorie AUC, et \inserted{définissent} un 2-$\unsure{\otimes|O}$-Champ AUC. Il s'envoie dans la 1-\inserted{$\otimes$-}cat. AUC (ou 1-$\otimes$-champ AUC) des $(M,N,\sigma\colon M\longrightarrow{}_2N)$, mais si $\Lambda,\Lambda'\colon\mathscr P\longrightarrow\mathscr P'$ définissent le $=$
\[
(M,N,\sigma)\xrightarrow{\lambda}(M',N',\sigma'),
\]
\deleted{\unsure{on ne peut plus}}\inserted{(p. ex. l'identité)}, on
% archive sheet number: 13
\scanpage{15}
auquel on se ramène) on ne peut plus affirmer que $\Lambda$ et $\Lambda'$ sont isomorphes \deleted{\unsure{ce}} (ni $=$ loc. isom.); c'est \inserted{tjrs vrai (pr un \deleted{$\unsure{\mathscr P|P},\mathscr P'\ (M,N,\sigma)\to(M',N',\sigma')$ fixés)} $\Lambda\colon\mathscr P\to\mathscr P'$ fixé)} \deleted{\unsure{sauf}} \unsure{sss|ssi|si} $\operatorname{Ext}^1(S;M,N')=0$ (resp. $\underline{\operatorname{Ext}}^1(M,N')=0$). \deleted{Mais si on se donne}\inserted{On ne peut} affirmer non plus que tt hom $\lambda$ de $(M,N,\sigma)$ ds $(M',N',\sigma')$ provienne d'un $\Lambda$ \inserted{($=$ loc.)} \inserted{[c'est OK ssi $\operatorname{Ext}^{\unsure{2|1}}(S;\unsure{M},N')=0$ resp. $\underline{\operatorname{Ext}}^2(M,N')=0$]} (ni que tt $(M,N,\sigma)$ provienne d'un $\mathscr P$ -- cf. plus bas --)
\margin{On peut prendre pour $\lambda$ l'identité}
mais pour $\mathscr P,\mathscr P',\lambda$ fixé, le\deleted{\unsure{s}} \deleted{cat. des} champ des $\otimes$-hom AUC $\mathscr P\to\mathscr P'$ correspondant est un 1-\inserted{ps.}torseur sous le champ de Picard \deleted{des}\inserted{strict} d'invariants $\underline{\operatorname{Ext}}^1(M,N'),\underline{\operatorname{Hom}}(M,N')$ obtenu par tronçage de $R\underline{\operatorname{Hom}}(M,N')$; l'obstruction à trouver un tel \deleted{\unsure{hom}}\inserted{$\Lambda$} est dans $\operatorname{Ext}^2(S,M,N')$ [et à le trouver loc est dans $H^0(S,\underline{\operatorname{Ext}}^2(M,N'))$].

Tout ceci provient du jeu habituel à la Baer, et du fait (qui en résulte) que pour $M,N,\sigma$ \deleted{\unsure{$(M',N',\sigma')$}} fixés, les champs de Picard d'invariants
% archive sheet number: 14
\scanpage{16}
\mspage{5}
$M,N,\sigma$ forment une 2-catégorie \deleted{2-équivalente à celle}\inserted{qui est un 2-ps-tors sous la catégorie de Picard} des champs de Picard \uline{stricts} d'inv. $M,N$ (et 0 donc) -- la ``différence'' de \deleted{deux} $\unsure{\mathscr P|P}$ et $\mathscr P'$ (épinglés par $M,N,\sigma$) est donc une \unsure{Cat|cat.} de Picard stricte d'inv. $M,N$ -- d'invariant $\in\operatorname{Ext}^2(S;M,N)$)

\noindent ----

\deleted{\unsure{§}} Il reste encore à voir à quelles conditions un triple $(M,N,\sigma)$ provient d'un $\mathscr P$. \deleted{(L'obstruction serait-elle dans $\operatorname{Ext}^3(S;M,N)$ ?)}
% slanted
\margin{cette condition est satisfaite ssi $\mathscr P$ (d'invariant $(M,N,\sigma)$) admet une deuxième \deleted{\unsure{contrainte de comm.}} \unsure{Comm.|comm.} qui en fasse un champ de Picard \uuline{strict}.}
(C'est OK si $\sigma$ est la restriction d'une application \inserted{bilin} antisymétrique $M\times M\longrightarrow N$ $\varphi(x,y)=-\varphi(y,x)$ -- c'est le cas si $M_2$ \deleted{\unsure{$M$ tel $M_2$ soit un module projectif sur}} $\unsure{)}$ est un \unsure{$\mathbb F_2$-}faisceau ``libre'', p.ex. si $M$ est ``libre'' sur $\mathbb Z$. Cela montre déjà que $M$ est quotient d'un $L$ tel que le composé $L\overset{p}{\unsure{\relbar|\longrightarrow}}M\xrightarrow{\sigma}{}_2N$ est réalisable par une cat. de Picard \inserted{$\mathscr P$} (déterminée \unsure{près|par|pour} un élément de $\operatorname{Ext}^2(S;L,N)$) La ``restriction'' de cette $\mathscr P$ à $(R,N)$ ($R=\operatorname{Ker}p$) définit une classe ds $\operatorname{Ext}^2(S;R,N)$ déterminée mod images d'éléments de $\operatorname{Ext}^{\unsure{2}}(S;L,N)$ donc son image \inserted{$\operatorname{ob}(\sigma)$} ds $\operatorname{Ext}^3(S;M,N)$ est can. déterminée, et elle est nulle ssi l'on
% archive sheet number: 15
